\documentclass[11pt]{amsart}
\usepackage[margin=1.1in]{geometry}
\usepackage{amsmath,amssymb,amsthm,mathtools}
\usepackage{enumitem}
\usepackage{microtype}
\usepackage{hyperref}
\usepackage[nameinlink,capitalize]{cleveref}
\hypersetup{colorlinks=true,linkcolor=blue,citecolor=blue,urlcolor=blue}

\theoremstyle{plain}
\newtheorem{theorem}{Theorem}[section]
\newtheorem{lemma}[theorem]{Lemma}
\newtheorem{corollary}[theorem]{Corollary}
\newtheorem{assumption}[theorem]{Assumption}
\theoremstyle{definition}
\newtheorem{definition}[theorem]{Definition}
\theoremstyle{remark}
\newtheorem{remark}[theorem]{Remark}

\title[Local Caccioppoli and Rough-Core Exclusion]{Local Caccioppoli and Rough-Core Exclusion}
\author{}
\date{}

\begin{document}

\begin{abstract}
We present a fully explicit local Caccioppoli estimate for renormalized Navier--Stokes fields. All terms produced by the localized energy inequality are estimated separately, and all absorptions are written explicitly. The argument uses standard PDE notation and the usual integrability classes of suitable weak solutions.

In particular, we derive local windowed $H^1$ control on compact space--time cylinders from local compact Caffarelli--Kohn--Nirenberg quantities. The result gives the rough-core alternative as the negation of local compact-CKN control.

This paper is written in the classical setting of suitable weak solutions (Leray~\cite{Leray1934}, Caffarelli--Kohn--Nirenberg~\cite{CaffarelliKohnNirenberg1982}, Scheffer~\cite{Scheffer1976}, Lin~\cite{Lin1998}).
\end{abstract}

\maketitle

\section{Imported inputs and notation}

This paper uses the renormalized local energy inequality obtained in the preceding
renormalized-chart analysis as an imported input. It does not rederive the
renormalized suitable structure or the pressure reconstruction. All estimates
below are conditional on the following local hypotheses on the compact
renormalized cylinders under consideration.

\begin{assumption}[Renormalized local energy inequality]\label{ass:lei}
Let $(V,P,a,b)$ be defined on a compact renormalized cylinder
$Q_{2R,J}:=B_{2R}\times J$, where $J\subset\mathbb R$ is a bounded interval.
Assume
\begin{equation}\label{eq:renorm-ns}
\partial_\tau V+(V\cdot\nabla)V+\nabla P-\nu\Delta V
= a(\tau)V+a(\tau)(y\cdot\nabla V)+b(\tau)\cdot\nabla V,
\qquad
\nabla\cdot V=0,
\end{equation}
with $\nu>0$, $\nabla\cdot V=0$, and $a,b\in L^\infty(J)$. For every
nonnegative $\phi\in C_c^\infty(Q_{2R,J})$ and every $\tau_1<\tau_2$ in $J$,
\begin{equation}\label{eq:local-energy}
\begin{aligned}
&\int_{B_{2R}}\frac{|V(\cdot,\tau_2)|^2}{2}\phi(\cdot,\tau_2)
-\int_{B_{2R}}\frac{|V(\cdot,\tau_1)|^2}{2}\phi(\cdot,\tau_1)
\\
&\quad+\nu\int_{\tau_1}^{\tau_2}\!\int_{B_{2R}}\phi\,|\nabla V|^2
\le
\int_{\tau_1}^{\tau_2}\!\int_{B_{2R}}\frac{|V|^2}{2}(\partial_\tau\phi+\nu\Delta\phi)
\\
&\quad+\int_{\tau_1}^{\tau_2}\!\int_{B_{2R}}\frac{|V|^2}{2}V\cdot\nabla\phi
+\int_{\tau_1}^{\tau_2}\!\int_{B_{2R}}PV\cdot\nabla\phi
\\
&\quad+\int_{\tau_1}^{\tau_2}\!\int_{B_{2R}}
a\frac{|V|^2}{2}(-\phi-y\cdot\nabla\phi)
\\
&\quad-\int_{\tau_1}^{\tau_2}\!\int_{B_{2R}}\frac{|V|^2}{2}b\cdot\nabla\phi.
\end{aligned}
\end{equation}
The pressure is understood modulo functions of time, and the mean-normalized
pressure on compact balls belongs to the local class used below.
\end{assumption}

Throughout, for a space--time set $E\subset\mathbb R^3\times\mathbb R$, write
\[
\int_E f:=\int_E f(y,\tau)\,dy\,d\tau.
\]
For $r>0$ and a bounded interval $J\subset\mathbb R$, define
\[
Q_{r,J}:=B_r\times J,\qquad |Q_{r,J}|:=|B_r|\,|J|.
\]
For a ball $B\subset\mathbb R^3$, denote the spatial mean of pressure by
\[
(P)_B(\tau):=\frac1{|B|}\int_B P(y,\tau)\,dy.
\]
The local quantities are
\[
\mathcal A_J(r):=\operatorname*{ess\,sup}_{\tau\in J}\int_{B_r}|V(y,\tau)|^2\,dy,
\]
\[
\mathcal H_J(r):=\int_J\!\int_{B_r}|\nabla V|^2\,dy\,d\tau,
\]
\[
\mathcal C_J(r):=\int_J\!\int_{B_r}|V|^3\,dy\,d\tau,
\qquad
\mathcal D_J(r):=\int_J\!\int_{B_r}|P-(P)_{B_r}(\tau)|^{3/2}\,dy\,d\tau.
\]

\begin{lemma}[Finite-measure conversion]\label{lem:l2-from-l3}
For every bounded interval $J$ and every $r>0$,
\[
\int_J\!\int_{B_r}|V|^2
\le |Q_{r,J}|^{1/3}\mathcal C_J(r)^{2/3}
\le |Q_{r,J}|^{1/3}\bigl(1+\mathcal C_J(r)\bigr).
\]
\end{lemma}

\begin{proof}
H\"older's inequality on $Q_{r,J}$ gives
\[
\int_J\!\int_{B_r}|V|^2
\le |Q_{r,J}|^{1/3}\left(\int_J\!\int_{B_r}|V|^3\right)^{2/3}.
\]
The second inequality follows from $s^{2/3}\le 1+s$ for $s\ge0$.
\end{proof}

\begin{lemma}[Comparison of pressure normalizations]\label{lem:pressure-means}
Let $1\le p<\infty$ and $0<r\le R$. For almost every $\tau$ such that
$P(\tau)\in L^p(B_R)$,
\[
\|P(\tau)-(P)_{B_r}(\tau)\|_{L^p(B_r)}
\le 2\|P(\tau)-(P)_{B_R}(\tau)\|_{L^p(B_R)}.
\]
Consequently,
\[
\mathcal D_J(r)\le 2^{3/2}\mathcal D_J(R)
\]
whenever the right-hand side is finite.
\end{lemma}

\begin{proof}
Set $c=(P)_{B_R}(\tau)$. Then
\[
\|P-(P)_{B_r}\|_{L^p(B_r)}
\le \|P-c\|_{L^p(B_r)}+|B_r|^{1/p}|(P)_{B_r}-c|.
\]
Since
\[
|(P)_{B_r}-c|\le |B_r|^{-1}\int_{B_r}|P-c|
\le |B_r|^{-1/p}\|P-c\|_{L^p(B_r)},
\]
the displayed pointwise estimate follows. Raising to the power $3/2$ and
integrating in time gives the last assertion.
\end{proof}

\section{Compact-cylinder Caccioppoli estimate}

\begin{theorem}[Compact-cylinder Caccioppoli estimate]\label{thm:caccioppoli}
Let $R>0$, let $J\subset\mathbb R$ be a bounded open interval, and let
$J'\Subset J$ be a compact subinterval. Assume
\cref{ass:lei} on $Q_{2R,J}$ and assume
\[
\mathcal C_J(2R)+\mathcal D_J(2R)<\infty.
\]
Then there is a constant
\[
C=C\bigl(R,|J|,\nu,\|a\|_{L^\infty(J)},\|b\|_{L^\infty(J)},
\operatorname{dist}(J',\partial J)\bigr)
\]
such that
\[
\mathcal A_{J'}(R)+\mathcal H_{J'}(R)
\le C\Bigl(1+
\mathcal C_J(2R)+\mathcal D_J(2R)
\Bigr).
\]
The estimate is outer-to-inner: all right-hand quantities are measured on
$B_{2R}\times J$, while the conclusion is on $B_R\times J'$.
\end{theorem}

\begin{proof}
\textbf{Step 1: localized cut-off and the two energy inequalities.}
Set
\[
\delta:=\operatorname{dist}(J',\partial J)>0.
\]
Choose $\eta\in C_c^\infty(J)$ and $\zeta\in C_c^\infty(B_{2R})$ with
\[
0\le\eta\le1,\qquad
0\le\zeta\le1,\qquad
\eta\equiv1\ \text{on}\ J',\qquad
\zeta\equiv1\ \text{on}\ B_R,
\]
and
\[
\|\eta'\|_{L^\infty(J)}\le C_\eta\,\delta^{-1},
\qquad
\|\nabla\zeta\|_{L^\infty(B_{2R})}\le C_\zeta R^{-1},
\qquad
\|\Delta\zeta\|_{L^\infty(B_{2R})}\le C_\zeta R^{-2}.
\]
Set
\[
\phi(y,\tau):=\eta(\tau)\,\zeta(y)^2.
\]
Then $\phi\in C_c^\infty(Q_{2R,J})$, $0\le\phi\le1$, and
\[
\phi\equiv1\text{ on }B_R\times J'.
\]
Choose $\tau_-<\inf J'$ and $\tau_+>\sup J'$ in $J$ with
$\eta(\tau_-)=\eta(\tau_+)=0$. Applying \eqref{eq:local-energy} on
$(\tau_-,\tau_+)$ gives
\[
\nu\int_{\tau_-}^{\tau_+}\!\int_{B_{2R}}\phi\,|\nabla V|^2
\le |R_\tau|+|R_\Delta|+|R_{\mathrm{conv}}|+|R_{\mathrm{pres}}|+|R_{\mathrm{mod}}|,
\]
where
\[
R_\tau:=\int_J\!\int_{B_{2R}}\frac{|V|^2}{2}\,\partial_\tau\phi,\qquad
R_\Delta:=\nu\int_J\!\int_{B_{2R}}\frac{|V|^2}{2}\,\Delta\phi,
\]
\[
R_{\mathrm{conv}}:=\int_J\!\int_{B_{2R}}\bigl(\tfrac{|V|^2}{2}V\bigr)\cdot\nabla\phi,\qquad
R_{\mathrm{pres}}:=\int_J\!\int_{B_{2R}}\bigl(P-(P)_{B_{2R}}(\tau)\bigr)V\cdot\nabla\phi,
\]
\[
R_{\mathrm{mod}}:=\int_J\!\int_{B_{2R}}
a\frac{|V|^2}{2}(-\phi-y\cdot\nabla\phi)
-\int_J\!\int_{B_{2R}}\frac{|V|^2}{2}b\cdot\nabla\phi.
\]
For almost every $\sigma\in J'$, applying \eqref{eq:local-energy} on
$(\tau_-,\sigma)$ gives the same right-hand estimates and leaves the terminal
term $\frac12\int_{B_R}|V(\sigma)|^2$ on the left. Therefore the bounds below
control both $\mathcal H_{J'}(R)$ and $\mathcal A_{J'}(R)$.
Define
\[
L:=\int_J\!\int_{B_{2R}}|V|^2,\qquad
C_3:=\mathcal C_J(2R),\qquad D_{3/2}:=\mathcal D_J(2R).
\]
By \cref{lem:l2-from-l3},
\[
L\le |Q_{2R,J}|^{1/3}C_3^{2/3}\le C(R,|J|)(1+C_3).
\]

Since $\phi(y,\tau)=\eta(\tau)\zeta(y)^2$,
\[
\partial_\tau\phi=\eta'(\tau)\zeta(y)^2,
\qquad
\nabla\phi=2\eta(\tau)\zeta(y)\nabla\zeta(y),
\qquad
\Delta\phi=\eta(\tau)\Delta(\zeta(y)^2),
\]
hence
\[
\|\partial_\tau\phi\|_{L^\infty(Q_{2R,J})}\le C_\eta\delta^{-1},
\qquad
\|\nabla\phi\|_{L^\infty(Q_{2R,J})}\le 2C_\zeta R^{-1},
\qquad
\|\Delta\phi\|_{L^\infty(Q_{2R,J})}\le C_\zeta R^{-2}.
\]

\textbf{Step 2: derivative terms from the cut-off.}
\[
|R_\tau|
\le\frac12\|\partial_\tau\phi\|_{L^\infty(Q_{2R,J})}L
\le C\delta^{-1}(1+C_3),
\]
\[
|R_\Delta|
\le\frac{\nu}{2}\|\Delta\phi\|_{L^\infty(Q_{2R,J})}L
\le C\nu R^{-2}(1+C_3).
\]

\textbf{Step 3: convection term.}
Using $|\nabla\phi|\le CR^{-1}$,
\[
|R_{\mathrm{conv}}|
\le \frac12\|\nabla\phi\|_{L^\infty(Q_{2R,J})}\int_J\!\int_{B_{2R}}|V|^3
\le CR^{-1}C_3.
\]

\textbf{Step 4: pressure term.}
For each time $\tau$,
\[
\int_{B_{2R}}(P-(P)_{B_{2R}}(\tau))V\cdot \nabla\phi
=\int_{B_{2R}}PV\cdot \nabla\phi,
\]
with
\[
\int_{B_{2R}}(P)_{B_{2R}}(\tau)V\cdot \nabla\phi
=(P)_{B_{2R}}(\tau)\int_{B_{2R}}V\cdot \nabla\phi
=-(P)_{B_{2R}}(\tau)\int_{B_{2R}}\phi\,\nabla\cdot V=0
\]
(using $\phi\in C_c^\infty(B_{2R})$ and $\nabla\cdot V=0$).
Thus subtracting any time-dependent pressure constant is legitimate in the
pressure flux term.
Hence
\[
|R_{\mathrm{pres}}|
\le\|\nabla\phi\|_{L^\infty(Q_{2R,J})}
\int_J
\|P-(P)_{B_{2R}}(\tau)\|_{L_x^{3/2}(B_{2R})}
\|V(\tau,\cdot)\|_{L_x^3(B_{2R})}
\,d\tau
\le CR^{-1}D_{3/2}^{2/3}C_3^{1/3}.
\]
Young's inequality gives
\[
|R_{\mathrm{pres}}|\le C(R)(C_3+D_{3/2}).
\]

\textbf{Step 5: modulation terms.}
The imported local energy inequality already contains the modulation terms
after the spatial integrations by parts from the repaired-gauge equation:
\[
\int_{B_{2R}}(y\cdot\nabla V)\cdot V\,\phi
=\frac12\int_{B_{2R}}y\cdot\nabla(|V|^2)\phi
=-\frac12\int_{B_{2R}}|V|^2\nabla\cdot(y\phi),
\]
and
\[
\int_{B_{2R}}(b(\tau)\cdot\nabla V)\cdot V\,\phi
=\frac12\int_{B_{2R}}b(\tau)\cdot\nabla(|V|^2)\phi
=-\frac12\int_{B_{2R}}|V|^2b(\tau)\cdot\nabla\phi.
\]
Thus we estimate the imported terms directly. Split
\[
R_{\mathrm{mod}}=R_a+R_{ay}+R_b,
\]
where
\[
R_a:=-\int_J\!\int_{B_{2R}}a\,\frac{|V|^2}{2}\phi,
\quad
R_{ay}:=-\int_J\!\int_{B_{2R}}a\,\frac{|V|^2}{2}\,y\cdot\nabla\phi,
\quad
R_b:=-\int_J\!\int_{B_{2R}}\frac{|V|^2}{2}b\cdot\nabla\phi.
\]
\[
M_a:=\|a\|_{L^\infty(J)},\qquad M_b:=\|b\|_{L^\infty(J)}.
\]
For $R_a$,
\[
|R_a|
\le \frac12 M_aL
\le C M_a(1+C_3).
\]
For the scale-gradient part, $|y|\le2R$ on $B_{2R}$ and
$\|\nabla\phi\|_{L^\infty(Q_{2R,J})}\le CR^{-1}$, hence
\[
|y\cdot\nabla\phi|\le C.
\]
Therefore
\[
|R_{ay}|\le C M_aL\le C M_a(1+C_3).
\]
For $R_b$,
\[
|R_b|\le C M_bR^{-1}L\le C M_bR^{-1}(1+C_3).
\]

\textbf{Step 6: absorption and conclusion.}
Collecting estimates from Steps~2--5 yields
\[
\nu\int_{\tau_-}^{\tau_+}\!\int_{B_{2R}}\phi|\nabla V|^2
\le C\bigl(1+C_3+D_{3/2}\bigr),
\]
where $C$ depends only on
\[
R,|J|,\nu,\|a\|_{L^\infty(J)},\|b\|_{L^\infty(J)},\delta.
\]
Since $\phi\equiv1$ on $B_R\times J'$,
\[
\mathcal H_{J'}(R)
\le C\bigl(1+\mathcal C_J(2R)+\mathcal D_J(2R)\bigr).
\]
Repeating the same estimates with terminal time $\sigma\in J'$ gives
\[
\frac12\int_{B_R}|V(y,\sigma)|^2\,dy
\le C\bigl(1+\mathcal C_J(2R)+\mathcal D_J(2R)\bigr)
\]
for almost every $\sigma\in J'$. Taking the essential supremum over
$\sigma\in J'$ and renaming $C$ proves the theorem.
\end{proof}

\begin{definition}[Local compact-cylinder control]
Local Caffarelli--Kohn--Nirenberg control on $Q_{r,J}$ means
\[
\mathcal C_J(r)+\mathcal D_J(r)<\infty.
\]
Local windowed $H^1$ control on $Q_{r,J}$ means
\[
\mathcal H_J(r)<\infty.
\]
\end{definition}

\begin{corollary}[Outer control implies inner suitable control]\label{cor:outer-inner}
Let $R>0$, let $J\subset\mathbb R$ be a bounded open interval, and let
$J'\Subset J$ be a compact subinterval. Assume \cref{ass:lei} on $Q_{2R,J}$ and
$\mathcal C_J(2R)+\mathcal D_J(2R)<\infty$. Then
\[
\mathcal A_{J'}(R)+\mathcal H_{J'}(R)<\infty.
\]
More quantitatively, the bound is the one in \cref{thm:caccioppoli}. If a
later argument uses pressure means on a smaller ball, \cref{lem:pressure-means}
converts those normalizations to the outer-ball normalization used here.
\end{corollary}

\begin{proof}
This is \cref{thm:caccioppoli}. The final sentence follows from
\cref{lem:pressure-means}.
\end{proof}

\begin{corollary}[Rough-core exclusion on compact windows]\label{cor:rough-core}
Let $R>0$, let $J\subset\mathbb R$ be a bounded open interval, and let
$J'\Subset J$ be a compact subinterval. Assume \cref{ass:lei} on $Q_{2R,J}$. If
\[
\mathcal C_J(2R)+\mathcal D_J(2R)<\infty,
\]
then
\[
\mathcal H_{J'}(R)<\infty.
\]
Equivalently, if $\mathcal H_{J'}(R)=+\infty$ for an inner cylinder
$Q_{R,J'}$, then for every enclosing cylinder $Q_{2R,J}$ with
$J'\Subset J$ on which the imported local energy inequality and bounded
modulation hypotheses hold,
\[
\mathcal C_J(2R)=+\infty
\qquad\text{or}\qquad
\mathcal D_J(2R)=+\infty.
\]
\end{corollary}

\begin{proof}
The direct implication is the gradient part of \cref{thm:caccioppoli}. The
last assertion is its contrapositive with the same outer-to-inner geometry:
finite outer $\mathcal C+\mathcal D$ on $Q_{2R,J}$ forces finite inner
$\mathcal H$ on $Q_{R,J'}$.
\end{proof}

\begin{remark}[Meaning of rough core]
In later applications, a rough core means failure of local windowed $H^1$
control on compact inner renormalized cylinders while the modulation
coefficients are bounded. \Cref{cor:rough-core} shows that such a failure is
not independent of compact-cylinder Caffarelli--Kohn--Nirenberg control: it can
occur only if the velocity-pressure package $\mathcal C+\mathcal D$ fails on
every relevant larger enclosing compact cylinder.
\end{remark}

\section{Inputs exported to later papers}

Later arguments may use the following consequences of this paper.
\begin{enumerate}[label=\textup{(\arabic*)}]
\item On every compact renormalized window satisfying the imported local energy
inequality and bounded modulation, finite $\mathcal C+\mathcal D$ on
$Q_{2R,J}$ implies finite $\mathcal A+\mathcal H$ on $Q_{R,J'}$ whenever
$J'\Subset J$.
\item The pressure normalization may be taken on any comparable ball, after
using \cref{lem:pressure-means}.
\item Failure of local windowed $H^1$ control on an inner compact cylinder
forces failure of the compact-cylinder $\mathcal C+\mathcal D$ package on each
larger enclosing cylinder to which \cref{thm:caccioppoli} applies.
\end{enumerate}

\bibliographystyle{alpha}
\bibliography{paper4_local_regularity_rough_core_exclusion}

\end{document}
